\documentclass[12pt]{article}
\usepackage[a4paper,margin=1in]{geometry}
\usepackage[T1]{fontenc}
\usepackage{lmodern,microtype}
\usepackage{amsmath,amssymb,amsthm,mathtools}
\usepackage{booktabs,array,enumitem,xcolor}
\usepackage[numbers,sort&compress]{natbib}
\usepackage[colorlinks=true,linkcolor=blue!55!black,citecolor=blue!55!black,urlcolor=blue!55!black]{hyperref}
\linespread{1.07}

\newtheorem{theorem}{Theorem}[section]
\newtheorem{proposition}[theorem]{Proposition}
\newtheorem{corollary}[theorem]{Corollary}
\theoremstyle{remark}
\newtheorem{remark}[theorem]{Remark}

\title{Source-Channel Determinant and Trace Factorization\\
\large Separating Spectral Newton Sums from Residue-Weighted Physical Moments}
\author{Liang Wang}
\date{July 2026}

\begin{document}
\maketitle

\begin{abstract}
The canonical source-channel packet carries two related but different data
sets.  Its spectral determinant and Newton traces are
\[
 D_E(z)=\det(zI-K)=\prod_j(z-\lambda_j),
 \qquad
 \operatorname{tr}K^q=\sum_j\lambda_j^q.
\]
The physical source-observation moments retain the channel residues:
\[
 h_q=c^*K^qb=\sum_jr_j\lambda_j^q.
\]
Conflating these ledgers would erase the input-output weights needed by any
later trace formula.

This paper gives an exact determinant encoding of the weighted transfer:
\[
 \frac{\det(zI-(K+bc^*))}{\det(zI-K)}
 =1-c^*(zI-K)^{-1}b.
\]
Thus one finite pair of characteristic determinants contains both the pole
set and the residue-weighted response.  A 240-case complex nonnormal audit
verifies the determinant lemma, modal weighted moments, and Newton traces
with zero failures.

For the physical length-four windows, the temporal determinant and traces
converge to the exact quartet ledger despite canonical frame norm products
as large as $950$.  At the latest starts the relative determinant error is
below $10^{-4}$ and the maximum relative trace error through power eight is
below $8\times10^{-4}$.

These are finite physical channel identities.  They do not provide
von Mangoldt weights, a prime-power orbit formula, a zeta determinant, or a
Hilbert--P\'olya operator.
\end{abstract}

\section{Two ledgers that must not be conflated}

Let $K\in\mathbb C^{r\times r}$ be the exact canonical packet of RH-196, and
let $b,c\in\mathbb C^r$ be source and observation coordinates.  In a
residue-labeled simple-mode basis one may take
\begin{equation}\label{eq:modal-data}
 K=\operatorname{diag}(\lambda_1,\ldots,\lambda_r),
 \qquad b=(1,\ldots,1)^T,
 \qquad \overline{c_j}=r_j.
\end{equation}

The spectral ledger ignores source and observation:
\begin{equation}\label{eq:spectral-ledger}
 D_E(z)=\det(zI-K),
 \qquad s_q=\operatorname{tr}K^q.
\end{equation}
The physical transfer ledger is
\begin{equation}\label{eq:transfer-ledger}
 g_E(z)=c^*(zI-K)^{-1}b,
 \qquad h_q=c^*K^qb.
\end{equation}
Both are similarity invariant if $b,c$ transform with the state coordinates,
but they answer different questions.

\section{Spectral determinant and Newton traces}

\begin{proposition}[Finite channel determinant]\label{prop:determinant}
For the selected simple modes,
\begin{equation}\label{eq:determinant}
 D_E(z)=\prod_{j=1}^r(z-\lambda_j).
\end{equation}
Its logarithmic derivative is
\begin{equation}\label{eq:log-derivative}
 \frac{D_E'(z)}{D_E(z)}
 =\operatorname{tr}(zI-K)^{-1}
 =\sum_{j=1}^r\frac{1}{z-\lambda_j}.
\end{equation}
\end{proposition}

The logarithmic derivative assigns unit multiplicity to each channel pole.
It does not recover the physical residues $r_j$ unless all happen to equal
one.

For $|z|>\rho(K)$,
\begin{equation}\label{eq:trace-expansion}
 \operatorname{tr}(zI-K)^{-1}
 =\sum_{q\ge0}\frac{\operatorname{tr}K^q}{z^{q+1}}.
\end{equation}
Thus the Newton sums are the unweighted resolvent coefficients.

\section{Residue-weighted moments}

The source-observation transfer has expansion
\begin{equation}\label{eq:weighted-expansion}
 g_E(z)=\sum_{q\ge0}\frac{h_q}{z^{q+1}},
 \qquad h_q=c^*K^qb.
\end{equation}

\begin{theorem}[Modal moment identity]\label{thm:moments}
In the simple residue-labeled basis,
\begin{equation}\label{eq:modal-moments}
 \boxed{h_q=\sum_{j=1}^r r_j\lambda_j^q.}
\end{equation}
\end{theorem}

\begin{proof}
Insert \eqref{eq:modal-data} into $c^*K^qb$.
\end{proof}

The distinction $s_q$ versus $h_q$ is the finite prototype of a later trace
problem.  A prime-power trace formula would require specific arithmetic
weights analogous to von Mangoldt coefficients.  The physical residues in
this paper are not identified with those weights.

\section{A determinant ratio for the transfer}

The matrix determinant lemma gives the exact bridge.

\begin{theorem}[Rank-one feedback factorization]\label{thm:feedback}
For $z\notin\sigma(K)$,
\begin{equation}\label{eq:feedback}
 \boxed{
 \frac{\det(zI-(K+bc^*))}{\det(zI-K)}
 =1-c^*(zI-K)^{-1}b.
 }
\end{equation}
\end{theorem}

\begin{proof}
Factor
\[
 zI-(K+bc^*)=(zI-K)
 \left[I-(zI-K)^{-1}bc^*\right]
\]
and use $\det(I-uv^*)=1-v^*u$.
\end{proof}

Consequently the scalar transfer is encoded by a pair of finite spectral
determinants:
\begin{equation}\label{eq:transfer-pair}
 g_E(z)=1-\frac{D_{E,\rm fb}(z)}{D_E(z)}.
\end{equation}
Zeros of the numerator ratio reflect the feedback realization; poles are
the uncanceled channel eigenvalues.  This is a standard system-theoretic
identity, here attached to the physical source-observation packet
\citep{Kailath1980}.

\section{Gauge covariance}

Under an invertible state gauge
\begin{equation}\label{eq:gauge}
 K'=G^{-1}KG,
 \qquad b'=G^{-1}b,
 \qquad c'=G^*c,
\end{equation}
one has
\begin{equation}\label{eq:gauge-invariance}
 D_E'=D_E,
 \quad \operatorname{tr}(K')^q=\operatorname{tr}K^q,
 \quad (c')^*(zI-K')^{-1}b'=c^*(zI-K)^{-1}b.
\end{equation}
The determinant ratio is therefore invariant.  This explains why large
balanced frame norms need not spoil the finite determinant/trace ledger:
coordinate amplification cancels in exact similarity invariants.

\section{Complex identity audit}

The implementation uses dimensions 2--9 and thirty complex random systems
per dimension, for 240 cases.  It verifies:
\begin{enumerate}[leftmargin=2em]
\item the rank-one determinant ratio at an off-spectrum point;
\item equality of direct moments $c^*K^qb$ and modal residue sums through
 order seven;
\item equality of power traces and modal Newton sums through power eight.
\end{enumerate}
All cases pass the $10^{-8}$ tolerance.  The modal audit uses independently
computed left/right eigenvectors and therefore tests residue normalization as
well as the determinant lemma.

\section{Physical temporal transport}

RH-194 stores the temporal compressed matrix and its matched physical quartet
ledger for every accepted window \citep{WangRH194}.  Across all twelve
windows, the maximum relative determinant error is
$8.56\times10^{-4}$ and the maximum relative power-trace error through order
eight is $4.12\times10^{-3}$.

At the latest starts:
\begin{center}
\begin{tabular}{lrr}
\toprule
side & relative determinant error & maximum trace error\\
\midrule
left & $8.91\times10^{-5}$ & $7.66\times10^{-4}$\\
right & $4.47\times10^{-5}$ & $3.97\times10^{-4}$\\
\bottomrule
\end{tabular}
\end{center}
The determinant error is not monotone at every start, but both late values
are below $10^{-4}$.  The power-trace errors exhibit the same overall decay
as the subspace gaps in RH-198.

\section{What the finite determinant accomplishes}

The result supplies a strict finite object:
\begin{equation*}
 (D_E,D_{E,\rm fb})
 \quad\leftrightarrow\quad
 (\lambda_j,r_j)_{j=1}^4.
\end{equation*}
It retains both spectral locations and source-observation weights, is
independent of balanced gauge, and is approximated by the temporal packet at
the audited anchor.

This is precisely the kind of bookkeeping needed before any infinite
spectral determinant can be proposed.  It also prevents a common
overstatement: matching eigenvalues or unweighted traces is not enough to
derive an arithmetic explicit formula.

\section{Boundaries and next step}

The determinant is a degree-four polynomial per side.  It is not an entire
function with zeta growth, has no $T\log T$ zero count, and contains no
demonstrated prime-power weights.  The residues are physical transfer
weights, not von Mangoldt coefficients.

The next problem is selection: why does the finest audited model produce a
four-mode nonreal edge packet while the predeclared length-three attempts at
coarser scales fail?  RH-200 studies the conjugate-pair parity and edge-gap
geometry behind the quartet.  Cross-scale transport remains the central
open leaf of Gate A.

\bibliographystyle{plainnat}
\bibliography{references}
\end{document}
